\section{Affine angular profiles with every exact section decision}
\label{sec:v65-affine-caustic}
The preceding simple profile can be extended without assuming that
only one clearance seam is active, or that the critical point avoids a
section boundary. Retaining affine, rather than only homogeneous,
first-order margins is important: it permits the physical guard width
to be fixed while a small positive margin at the center varies.
We obtain a local formula for the entire vector of exact-label densities.
All estimates in this section are at a fixed word and finite count.

\subsection{A width-uniform comparison for smooth guards}
Write $\Theta$ for the original nondecreasing smooth guard, zero on
$(-\infty,1]$ and one on $[2,\infty)$. Let $\omega=\omega_\theta$.

\begin{lemma}[Affine thresholds and bounded-variation guards]
\label{lem:v65-affine-guard}
Suppose
\[
 f(r\omega)=\beta+r a\cdot\omega+E_r(\theta),\qquad
 a\ne0,\quad |E_r|\le Lr^2.
\]
Put $f_r^{\rm a}(\theta)=\beta+r a\cdot\omega$. Uniformly in
$\beta\in\R$ and every $s>0$,
\begin{align}\label{eq:v65-affine-guard}
 \int_0^{2\pi}
 |\Theta(f(r\omega)/s)-\Theta(f_r^{\rm a}(\theta)/s)|\,d\theta
      &\le C_a\sqrt{Lr},\notag\\
 |\{\theta:\one_{f(r\omega)>0}\ne\one_{f_r^{\rm a}(\theta)>0}\}|
      &\le C_a\sqrt{Lr}.
\end{align}
The same bound, summed over the scalar inputs, holds for finite
products and minima of such guard values. It also bounds the measure
on which a finite Boolean formula in the sign bits changes.
\end{lemma}
\begin{proof}
For every real threshold $t$, disagreement between $f>t$ and
$f_r^{\rm a}>t$ implies
$|a\cdot\omega-(t-\beta)/r|\le Lr$.
The circle estimate in the proof of
Lemma~\ref{lem:v65-angular-threshold} is uniform in this translated
threshold. This proves the sign assertion and a uniform bound for
all threshold disagreements. Since $\Theta$ has total variation one,
its Stieltjes representation is
$\Theta(x)=\int\one_{\{x>v\}}\,d\Theta(v)$, with immaterial
endpoint conventions. Tonelli and the threshold bound give the first
inequality, including when $s/r$ tends to zero or infinity.
For factors in $[0,1]$, the difference of products, or of minima,
is at most the sum of the individual differences. A Boolean formula
can change only if some input bit changes. Apply the union bound to
these input events, not to the number of output labels.
\end{proof}

\subsection{The finite physical margin description}
Let $z$ be a normal-to-normal critical point in the physical-side
closure of a selected word. Assume every selected incidence is positive;
there may be several simultaneous competing-hit seams and arbitrary
section-boundary coincidences. Use the smooth endpoint continuation
only to choose the Morse coordinates $y$ of the preceding section.
Choose $\chi_z>0$ and a sufficiently small Morse disk on which all
selected incidences exceed $\chi_z$. Fix
\begin{equation}\label{eq:v65-affine-width-range}
 0<\varepsilon<\min\{\chi_z/4,s_0/8\}.
\end{equation}
Every incidence guard is then one on the whole disk. All physical
defect source in this disk is therefore the sum of the original
first-clearance pieces, with their original ordering unchanged.

There is a finite list $g_i$ of smooth near-clearance margins, grouped
by their original flight $j(i)$, and a finite list $q_l$ of smooth
section-edge margins, with the following properties. A margin $g_i$
is the distance to its selected nonincident disk minus $R$, read at
image collision $j(i)+1$, with the perpendicular foot inside that
flight. Candidates which are uniformly too far to affect a guard in
\eqref{eq:v65-affine-width-range} are omitted. Each flight guard is
then
\begin{equation}\label{eq:v65-guard-minimum}
 U_j(y)=\min_{i:j(i)=j}
           \Theta(g_i(\Phi(y))/(\varepsilon e'_j)),
               \qquad \min\varnothing=1.
\end{equation}
The full physical guard on the disk is $\prod_jU_j$. The margins
$q_l$ are the rectangle edge functions for each original section
membership test, composed with its actual selected collision branch.
The section decisions at $0,\ldots,m-1$ define the occupation, and
terminal membership at $m$ is imposed separately.

\begin{lemma}[Nonzero scalar gradients and exact Boolean labels]
\label{lem:v65-margin-description}
The lists may be chosen so that
$D(g_i\circ\Phi)(0)\ne0$ and $D(q_l\circ\Phi)(0)\ne0$ for every
retained margin. The physical word and the original events
$\Pi_{n,k,m,R}$ are described on the disk by a finite Boolean formula
in their signs, with the fixed inactive conditions retained as constants.
Let its indicator be $B_{n,k}$. At every input sign pattern these
indicators can be defined by the original visit-count rule so that
\begin{equation}\label{eq:v65-boolean-partition}
 \sum_{n=1}^m\sum_k B_{n,k}\le1.
\end{equation}
The initial and terminal section tests are included in $B_{n,k}$;
no normalization is encoded in it.
\end{lemma}
\begin{proof}
A disk whose clearance can be below $s_0$ has its foot strictly
inside the flight, since the endpoint distances to a nonincident
disk are at least $\ell_0>s_0$. The finite-horizon candidate list
is fixed at each flight. On its near-clearance envelope,
$\partial_\varphi Z_d=L_d\ge\ell_*>0$ by
Lemma~\ref{lem:v49-physical-envelope}. Its sign is fixed locally
when $|Z_d|$ is near $R>0$. Recentring at the image collision and
passing to endpoint and Morse coordinates are local diffeomorphisms.
This proves the nonzero derivative for every needed $g_i$.
Candidates with clearance at least $s_0$ at $z$ remain above
$s_0/2$ after reducing the disk, and are inactive for
\eqref{eq:v65-affine-width-range}. Candidates with a foot outside
the segment likewise cannot enter the near-clearance range on a
sufficiently small disk. These observations justify the finite list
and \eqref{eq:v65-guard-minimum}; monotonicity of $\Theta$ allows it
to commute with the minimum of the actual clearance distances.

In the fixed section charts an edge function has nonzero coordinate
gradient. Its composition with a regular collision branch still has
nonzero gradient. Rectangle intersections and unions, followed by the
finite visit-count rule, give the claimed Boolean description. First-hit
exclusions on the selected simple-root word are exactly the relevant
positive clearance signs near the disk; other exclusions are strict.
All quantities are evaluated on that selected branch before imposing
its physical-side inequalities.

For any sign pattern, first evaluate each collision's section bit by
the same rectangle formula. If either endpoint is outside the section,
assign no exact return label. Otherwise the sum of bits at
$0,\ldots,m-1$ determines a unique $n$; the selected word determines
$k$. Assign precisely that label when the physical sign conditions
also hold. This construction is meaningful even for an affine sign
pattern not realized by an actual trajectory, and proves
\eqref{eq:v65-boolean-partition}. It assigns no measure to such a
pattern.
\end{proof}

Put $r=\sqrt{2h}$ and retain the constant terms in
\begin{align}\label{eq:v65-affine-margins}
 G_{i,r}(\theta)&=g_i(z)+rD(g_i\circ\Phi)(0)\cdot\omega_\theta,
                                      \notag\\
 Q_{l,r}(\theta)&=q_l(z)+rD(q_l\circ\Phi)(0)\cdot\omega_\theta.
\end{align}
Let $B_{n,k}^{\rm a}(r,\theta)$ be the same Boolean formula evaluated
on these affine signs. Define
\begin{align}\label{eq:v65-affine-profile}
 U_j^{\rm a}(\varepsilon,r,\theta)
   &=\min_{i:j(i)=j}\Theta(G_{i,r}(\theta)/(\varepsilon e'_j)),
                                      \notag\\
 \mathcal A_{z,n,k}(\varepsilon,h)
   &=\int_0^{2\pi}B_{n,k}^{\rm a}(r,\theta)
           \left(1-\prod_jU_j^{\rm a}(\varepsilon,r,\theta)\right)d\theta.
\end{align}
In particular $0\le\sum_{n,k}\mathcal A_{z,n,k}\le2\pi$.
These are explicit one-dimensional angular integrals of the physical
first jets and the inherited guard. They are not billiard probabilities.

\begin{theorem}[The full local exact-label caustic profile]
\label{thm:v65-affine-label-profile}
Let $b_{z,n,k}^{\varepsilon,\mathrm{clr}}$ be the original
first-clearance source restricted to the Morse disk at $z$, with
its exact labels. For suitable $h_z>0$ and $C_z<\infty$, uniformly
for \eqref{eq:v65-affine-width-range} and almost every $0<h<h_z$,
\begin{equation}\label{eq:v65-affine-label-error}
 \sum_{n=1}^m\sum_k
 \left|b_{z,n,k}^{\varepsilon,\mathrm{clr}}(t_z+h)
       -\frac{J_z}{2\pi}\mathcal A_{z,n,k}(\varepsilon,h)\right|
                         \le C_zJ_z h^{1/4}.
\end{equation}
No clearance or section margin at $z$ is required to be positive.
There is no uniform lower bound assumed on distinct margin gradients
or on angles between them. The theorem requires their individual
nonvanishing and keeps the resulting local constant $C_z$.
There is no factor for the number of output labels in this estimate.
More explicitly, take the Morse radius at most one, put
$W_z=\|D\widetilde w\|_\infty/\widetilde w(0)$, and for every scalar
margin $f$ in the two lists put
$A_f=|D(f\circ\Phi)(0)|>0$ and
$L_f=\tfrac12\|D^2(f\circ\Phi)\|_\infty$ on the disk.
An absolute constant $C$ permits the choice
\begin{equation}\label{eq:v65-jet-budget}
 C_z=C\left(1+W_z+\sum_f\sqrt{L_f/A_f}\right).
\end{equation}
Thus the error uses explicitly stated finite-jet data, not an
unrecorded separation exponent.
\end{theorem}
\begin{proof}
Use polar coarea as in \eqref{eq:v65-polar-source}, with the exact
integrand
\[
 \widetilde w(r\omega)B_{n,k}(\Phi(r\omega))
                 \left(1-\prod_jU_j(r\omega)\right).
\]
Taylor remainder bounds for the finite scalar lists are $O_z(r^2)$.
Lemma~\ref{lem:v65-affine-guard}, summed over these lists, bounds
the integral of the guard-product difference by $C_z r^{1/2}$,
uniformly in every width $\varepsilon e'_j$. It also bounds by
$C_zr^{1/2}$ the angular set on which any Boolean input sign differs
from its affine counterpart. Outside this set every exact-label
indicator agrees. On this set both the original indicator vector
and the affine one are either zero or a single unit vector by
\eqref{eq:v65-boolean-partition}. Their $\ell^1$ distance is
therefore at most two, regardless of the number of labels. This
bounds the summed indicator error by $2C_zr^{1/2}$.

Finally, $\widetilde w(r\omega)=J_z/(2\pi)+O_z(J_zr)$.
The summed output indicators are at most one, so its contribution
to the summed density error is $O_z(J_zr)$, not that quantity times
the number of labels. Combine the three estimates and use
$r^{1/2}=2^{1/4}h^{1/4}$. Every source restriction used in this
calculation is an original physical and section restriction. The
affine formulas approximate its angular integral; they do not alter
its itinerary or its pushforward measure. The circle threshold
estimate costs $C\sqrt{L_f/A_f}\,r^{1/2}$ for each scalar input.
The relative weight error costs $CW_zr$. Summing and using $r\le1$
gives \eqref{eq:v65-jet-budget} after an absolute enlargement of $C$.
\end{proof}

\begin{corollary}[Angular trace at a multiple seam]
\label{cor:v65-multiple-seam-trace}
Suppose at least one clearance margin vanishes at $z$ and keep
$\varepsilon$ fixed in \eqref{eq:v65-affine-width-range}. Replace
each zero margin in the Boolean label formula by the sign of its
linear part, and each nonzero margin by its constant sign. Let
$\mathcal C_{z,n,k}$ be the resulting angular sets. Then
\begin{equation}\label{eq:v65-cone-trace}
 \lim_{h\downarrow0}b_{z,n,k}^{\varepsilon,\mathrm{clr}}(t_z+h)
     =\frac{J_z}{2\pi}|\mathcal C_{z,n,k}|,
       \qquad \sum_{n,k}|\mathcal C_{z,n,k}|\le2\pi.
\end{equation}
The sets may have zero angular measure; no transverse intersection
between different active boundaries is assumed. For a sole active
clearance seam with strict section decisions this recovers $J_z/2$.
\end{corollary}
\begin{proof}
Each nonzero linear form vanishes at only two angular directions.
Off their finite union the affine signs converge to the stated signs.
On an allowed physical direction an active clearance margin is
$O(r)$; its guard tends to zero at fixed $\varepsilon$. Hence
$1-\prod_jU_j^{\rm a}$ tends to one on every allowed direction.
Dominated convergence in \eqref{eq:v65-affine-profile} and
\eqref{eq:v65-affine-label-error} prove the formula. The angular
partition bound follows from the same unique occupation rule.
Coincident linear forms can make some angular sets empty; the proof
never divides by a determinant formed from two such gradients.
\end{proof}

\paragraph{Uniformity, source coverage and the next dynamical estimate.}
The formula removes the single-seam and strict-section restrictions
from the local roof-critical calculation. It also avoids imposing
$\varepsilon$ below each nonzero clearance margin: the offsets
$g_i(z)$ in \eqref{eq:v65-affine-margins} retain such margins exactly
at first order. On a compact family of these coordinate neighborhoods,
with common derivative bounds and positive lower bounds for the stated
individual gradients and selected incidences, $h_z$ and $C_z$ can be
chosen uniformly by the proof. No such bounds are asserted for all
words as the collision count grows.

For disjoint original-source neighborhoods, the vector estimates may
be summed with their actual positive weights $J_z$. In particular a
local roof density is bounded by $(1+o_z(1))J_z$, and each $J_z$ is
the reversible square-root weight of
\eqref{eq:v65-reversible-weight}. Thus the remaining long-count task
is a weighted critical-cluster estimate together with uniform coverage
of the neighborhoods, not a conversion of their small mass into a
height estimate. This observation is a reduction for the described
roof-critical source, not a proof of either ordered physical height
limit. Regions with selected grazing contacts or noncritical-roof
rank degeneracy remain in the original positive clearance complement.
The complete incidence source, the canonical arithmetic kernel and
the fixed-band order of limits are unchanged.

\label{end:v65-new-mathematics}
